\section{The signal and its two noise scales}\label{sec:signal}

The signal must satisfy two competing requirements. Its higher cumulants
must be sparse enough for the Fisher estimate, while many coordinates
must share a block decomposition. Tensor independence supplies the latter
without spoiling the former.

Fix an integer $p\geq2$ and set $n=p+1$. Let $B_1,\ldots,B_p$ be free
variance-one semicirculars. In a separate tensor factor take a variable
$Y$ with uniform distribution on $p$ equally spaced points
$y_l=-1+2(l-1)/(p-1)$, $1\leq l\leq p$. Thus $Y$ commutes with every
$B_i$, but is not being taken free from them. Set
\[
 C=(B_1,\ldots,B_p,Y),\qquad V_a=C+\sqrt a S,\qquad
 L(a)=\frac12\sum_{i=1}^n\log(2\pi e(a+c_i)),
\]
where $c_i=\tau(C_i^2)$. Thus $c_i=1$ for $i\leq p$, and $C$ is
centered with diagonal covariance.
The benchmark $L(a)$ is the nonmicrostates entropy of the centered
semicircular family with coordinate variances $a+c_i$, and also the
normalized classical entropy of the corresponding independent GUE law.

The next construction provides a faithful normal tracial realization
together with uniformly bounded deterministic matrix models for $C$.
These models will later give the finite lower bound for $\chi(X)$.
Choose uniformly bounded deterministic matrix models
$B^{(q_k)}$ of the free semicircular $p$-tuple, using the convergence and
norm bounds in Section~\ref{sec:conventions}. Tensor them with $I_p$ and
take $I_{q_k}\otimes\operatorname{diag}(y_1,\ldots,y_p)$ for $Y$.
In dimensions $d_k=pq_k$ these give models of the required tensor law.
Adjoin an independent GUE $n$-tuple and select realizations satisfying
successively more moment tests and one fixed norm bound. Their joint
limiting moments give a positive tracial functional in the variables
$(C,S)$ with the asserted freeness. In its GNS construction, left and
right coordinate multiplication are bounded: the matrix inequality
$\tr_d(Q^*D_i^2Q)\leq R_i^2\tr_d(Q^*Q)$ passes to the limit.
The vacuum is cyclic for right polynomial multiplication, which commutes
with the weak closure of the left algebra, so the vacuum is separating
for that closure. Its vector state is therefore faithful and normal;
the trace identity extends by separate weak continuity. This supplies
the required tracial von Neumann algebra. Further independent noise
families can be included in the same construction. Faithfulness and the
semicircular distribution give $\|B_i\|=\|S_i\|=2$.

\begin{lemma}\label{lem:cumulants}
For $K=128$ and every $m\geq1$,
\[
 \|\kappa_m(C)\|_2\leq K^m p^{m/4}.
\]
\end{lemma}

\begin{proof}
There are at most $4^m$ noncrossing partitions and their M\"obius
coefficients have magnitude at most $4^m$, by the Catalan product formula
and the interval factorization of the noncrossing lattice
\cite[Proposition~9.4, Theorem~9.29, and Propositions~10.14--10.15]{NS}.
Each product of block moments has absolute value at most $2^m$, since
$\|B_i\|=2$ and $\|Y\|\leq1$. Thus each cumulant is at most $32^m$
in absolute value.

Suppose $b$ of the $m$ positions have $B$ indices. A nonzero product of
block moments in M\"obius inversion requires, after deleting the $Y$
positions in each block, a semicircular word with a noncrossing pairing
matching equal indices. Indeed, tensor independence factors each block
moment into its $B$-word moment and its $Y$-moment. The union of these
pairings is still noncrossing:
crossing pairs from distinct blocks would cross the containing partition.
For fixed $b$ positions, at most $4^b p^{b/2}$ labelings can therefore
contribute. Summing over positions bounds the number of potentially
nonzero entries by $(1+4\sqrt p)^m$. Consequently
\[
 \|\kappa_m(C)\|_2\leq32^m(1+4\sqrt p)^{m/2}
 \leq(32\sqrt5)^m p^{m/4}\leq K^m p^{m/4}.
\]
\end{proof}

By Proposition~\ref{prop:fisher} and a geometric series with ratio
$K^2\sqrt p/u<1/2$, for
$u>2K^2\sqrt p$ we have
\begin{equation}\label{eq:specialfisher}
 0\leq\Phi^*(V_u)-\sum_i\frac1{u+c_i}
 \leq 2K^6p^{3/2}u^{-4}.
\end{equation}
Adding independent free semicircular noise of variance $s$ to $V_a$
gives the law of $V_{a+s}$. Equation~\eqref{eq:chistar} therefore yields
\begin{equation}\label{eq:entropydeficit}
 \chi^*(V_a)=L(a)-\frac12\int_a^\infty
 \left[\Phi^*(V_u)-\sum_i\frac1{u+c_i}\right]du.
\end{equation}
All terms are finite for $a>0$: the integrand is measurable by
\eqref{eq:fishervariational}, nonnegative, and bounded above by
\[
 \frac n u-\sum_i\frac1{u+c_i}
 =\sum_i\frac{c_i}{u(u+c_i)}
 \leq\frac{\sum_i c_i}{u^2},
\]
which is integrable on $[a,\infty)$. In deriving
\eqref{eq:entropydeficit}, the covariance correction is
$\int_0^\infty[(1+s)^{-1}-(a+c_i+s)^{-1}]\,ds=\log(a+c_i)$;
thus no subtraction of divergent integrals is used.
In particular,
\begin{align}
 \chi^*(V_a)&\geq L(a)-\frac{K^6p^{3/2}}{3a^3}>L(a)-\tfrac12
 &&(a>2K^2\sqrt p),\label{eq:largeentropy}\\
 \chi^*(V_{h+t})-\chi^*(V_h)
 &=\frac12\int_0^t\Phi^*(V_{h+s})\,ds.
 \label{eq:freeflow}
\end{align}

We now fix $p$ sufficiently large and put $t=p^{3/4}$ so that
\begin{equation}\label{eq:parameters}
 \begin{gathered}
 t>2K^2\sqrt p,\qquad \frac{32t}{p}<\frac18,\\
 \frac{p}{256(2+t)}>
 2\log(1+128\sqrt{2+t})+\log2+8.
 \end{gathered}
\end{equation}
These inequalities are compatible because $p/t=p^{1/4}$ grows faster
than $\log p$. An explicit choice is $p=2^{64}$ and $t=2^{48}$:
the first threshold is $2^{47}$ and $32t/p=2^{-11}$.
For the last inequality, its left side exceeds $128$ since
$2+t<2^{49}$, whereas its right side is less than
$67\log2+8<75$ since $1+128\sqrt{2+t}<2^{33}$.
From this point on, $p$ and $t$ are fixed, and only matrix sizes will
tend to infinity.

For small $h>0$, partition the spectrum of $X_n=Y+\sqrt h S_n$ at the
midpoints between the $y_l$ and denote its spectral projections by $e_l(h)$.
Let $\Delta=2/(p-1)$ be the atom spacing. Since
$\|\sqrt h S_n\|=2\sqrt h$, the condition $h<\Delta^2/64$ places
the spectrum inside the disjoint $\Delta/4$-neighborhoods of the atoms,
with every threshold in a spectral gap. Choose continuous functions
agreeing with the interval indicators on these common neighborhoods.
Continuous functional calculus, or uniform polynomial approximation
there followed by norm continuity, gives $e_l(h)\to e_l(0)$ in norm.
At $h=0$ these projections commute with every $B_i$, have trace $1/p$,
and satisfy
\[
 \sum_{i=1}^p\sum_{l=1}^p\|e_l(0)B_i e_l(0)\|_2^2
 =\sum_{i=1}^p\|B_i\|_2^2=p.
\]
We may therefore fix $0<h<\min\{1,\Delta^2/64\}$ and abbreviate
$e_l=e_l(h)$ so that, for $X=V_h$,
\begin{equation}\label{eq:smallblocks}
 \tau(e_l)<\frac{3}{2p}\quad(1\leq l\leq p),\qquad
 \sum_{i=1}^p\sum_{l=1}^p\|e_lX_i e_l\|_2^2>\frac{3p}{4}.
\end{equation}

\begin{lemma}\label{lem:finitechi}
The tuple $X$ has $\chi(X)>-\infty$.
\end{lemma}

\begin{proof}
Use the deterministic models for $C$ above, in dimensions $d_k=pq_k$,
and add independent GUE noise of variance $h$.
If $M$ bounds the coordinate norms of those models, take
$R_0=M+8\sqrt h$. Lemma~\ref{lem:matrix-convergence} and
\eqref{eq:gue-norm-tail} show that these noised models lie in every
fixed moment neighborhood of $X$, within that cutoff, with probability
tending to one. Their density is bounded above by
$(d/(2\pi h))^{nd^2/2}$. If $q_d$ is the indicated probability, then
\[
 \Vol\Gamma_{R_0}(X;m,d,\varepsilon)
 \geq q_d(2\pi h/d)^{nd^2/2}.
\]
Since $q_d\to1$ along the model dimensions, the normalized limsup
logarithm is at least $\frac n2\log(2\pi h)$ for each $m,\varepsilon$.
Take the infimum and then the supremum over cutoffs. The limsup definition
makes this subsequence of matrix sizes sufficient.
\end{proof}
